\documentclass[10pt,a4paper]{article}
\usepackage[utf8]{inputenc}
\usepackage[margin=2.2cm]{geometry}
\usepackage[expansion=false]{microtype}
\usepackage{amsmath}
\usepackage{xcolor}
\usepackage{titlesec}
\usepackage[hidelinks]{hyperref}
\usepackage{url}

\titleformat{\section}{\normalsize\bfseries}{\thesection}{0.5em}{}
\titlespacing*{\section}{0pt}{8pt}{2pt}
\setlength{\parindent}{0pt}
\setlength{\parskip}{3pt}
\pagestyle{empty}

% visible placeholders, to be filled in before submission
\newcommand{\fillin}[1]{\textcolor{red}{[#1]}}

\begin{document}
\sloppy

{\large\bfseries Generalised Langford sequences as a CSP in MiniZinc}\\[1pt]
Symbolic Explainable AI 2026/27 \quad Group 20: António Delgado (106658), Antero Morgado (119213)\\
\textbf{CSPLib problem 024} (Langford's number problem)\\
Video: {\small\url{https://drive.google.com/drive/folders/1JhtKAuOsx3BNYTG6GkInrLacWVWrHRSs?usp=sharing}}

\section{Problem}
Two copies of each number $1..n$ must be placed in a sequence so that the two copies of $i$ are $i+1$ positions apart. Our starting point is the CSPLib model by Hakan Kjellerstrand (\texttt{langford2.mzn}), which only supports this classic case.

\section{Extension}
Our model (\texttt{langford\_gen.mzn}) generalises the problem to $s$ copies of each number, with consecutive copies of $i$ exactly $i+d$ positions apart ($s=2,d=1$ is the classic problem, $d=0$ the Nickerson variant). The decision variables are only the positions of the first copies, so the spacing always holds and the search works on $n$ variables. We added reversal symmetry breaking, which halves the search, and a parity constraint for the classic case, which rejects impossible instances ($n\not\equiv0,3\pmod4$) before search: $n=13$ takes 0.06\,s instead of 52\,s with the baseline. When there is no solution, \texttt{proj.py} also explains why in plain words. All results were checked against a brute-force enumerator; the benchmark is in \texttt{results.csv}.

\section{Running the code}
Requirements: Python 3.11.2, MiniZinc 2.10.1 (Gecode 6.4.0, Chuffed 0.14.0) on the \texttt{PATH}, and the \texttt{minizinc} Python package 0.10.0. Keep both \texttt{.mzn} files next to \texttt{proj.py}.
\vspace{-3pt}
{\footnotesize
\begin{verbatim}
python3 proj.py 8                       # first solution for n=8
python3 proj.py 9 -s 3 --all            # all solutions with 3 copies (d=1)
python3 proj.py 14 --parity             # impossible instance, with explanation
python3 proj.py --selftest              # compare with brute force
python3 proj.py --benchmark             # writes results.csv
\end{verbatim}}
\vspace{-8pt}

\section{Use of generative AI and web resources}
We used generative AI tools only to generate part of the code and to debug. The modelling and solving of the problem were done by us. For background we consulted \href{https://www.dialectrix.com/langford.html}{dialectrix.com/langford.html}, \href{https://www.dialectrix.com/langford/langford-biblio.html}{dialectrix.com/langford/langford-biblio.html}, \href{https://mathworld.wolfram.com/LangfordsProblem.html}{MathWorld} and \href{https://web.archive.org/web/19991002120421/http://www.lclark.edu/~miller/langford/MG.html}{lclark.edu/$\sim$miller/langford (archived)}.

\end{document}
